\documentclass{beamer}
\usepackage[utf8]{inputenc}

\usetheme{Madrid}
\usecolortheme{default}
\usepackage{amsmath,amssymb,amsthm}
\usepackage{txfonts}
\usepackage{tikz}
\usepackage{listings}
\usepackage{graphicx}
\usepackage{lmodern}

\setbeamertemplate{page number in head/foot}[totalframenumber]

\lstset{
    language=C,
    basicstyle=\ttfamily\small,
    keywordstyle=\color{blue},
    stringstyle=\color{orange},
    commentstyle=\color{green!60!black},
    numbers=left,
    numberstyle=\tiny\color{gray},
    breaklines=true,
    showstringspaces=false,
}

\title{5.5.25}
\author{Chiruvella Harshith Sharan - AI25BTECH11008}

\begin{document}

\frame{\titlepage}

%------------------------------------

\begin{frame}{Question}
Solve the system using vector algebra:
\[
\begin{aligned}
x + y + z &= 6 \\
2x + 3y + z &= 10 \\
x - y + 2z &= 5
\end{aligned}
\]
\end{frame}

%------------------------------------

\begin{frame}{Step 1: Matrix Form}
System can be written as
\[
A\vec{x} = \vec{b} \tag{1}
\]
where
\[
A = \begin{bmatrix}
1 & 1 & 1 \\
2 & 3 & 1 \\
1 & -1 & 2
\end{bmatrix}, 
\quad
\vec{x} = \begin{bmatrix} x \\ y \\ z \end{bmatrix}, 
\quad
\vec{b} = \begin{bmatrix} 6 \\ 10 \\ 5 \end{bmatrix}.
\]
\end{frame}

%------------------------------------

\begin{frame}{Step 2: Formula for Inverse}
The solution is
\[
\vec{x} = A^{-1}\vec{b} \tag{2}
\]
where
\[
A^{-1} = \frac{1}{\det(A)} \, \text{adj}(A).
\]
\end{frame}

%------------------------------------

\begin{frame}{Step 3: Computing Inverse}
\[
\det(A) = 7 \neq 0
\]

So inverse exists. After computation,
\[
A^{-1} =
\begin{bmatrix}
5/7 & 1/7 & -3/7 \\
1 & -1/7 & -1/7 \\
-2/7 & 2/7 & 3/7
\end{bmatrix} \tag{3}
\]
\end{frame}

%------------------------------------

\begin{frame}{Step 4: Solution}
\[
\vec{x} = A^{-1}\vec{b} =
\begin{bmatrix}
5/7 & 1/7 & -3/7 \\
1 & -1/7 & -1/7 \\
-2/7 & 2/7 & 3/7
\end{bmatrix}
\begin{bmatrix}
6 \\ 10 \\ 5
\end{bmatrix}
=
\begin{bmatrix}
1 \\ 2 \\ 3
\end{bmatrix} \tag{4}
\]
\end{frame}

%------------------------------------

\begin{frame}{Final Answer}
\[
\boxed{ \vec{x} = \begin{bmatrix} 1 \\ 2 \\ 3 \end{bmatrix} }
\]
\end{frame}

%------------------------------------

\begin{frame}[fragile]{Python Code}
\begin{lstlisting}[language=Python]
import matplotlib.pyplot as plt
from mpl_toolkits.mplot3d import Axes3D
import numpy as np

# Define column vectors of A
a1 = np.array([1,2,1])
a2 = np.array([2,0,2])
a3 = np.array([-3,-3,0])
b  = np.array([1,2,3])

# Scaled components from solution x = [2, 1/2, 2/3]
v1 = 2*a1
v2 = 0.5*a2
v3 = (2/3)*a3

# Step-by-step sum points
origin = np.array([0,0,0])
p1 = v1
p2 = v1+v2
p3 = v1+v2+v3  # should equal b


\end{lstlisting}
\end{frame}

\begin{frame}[fragile]{Python Code}
\begin{lstlisting}[language=Python]


fig = plt.figure(figsize=(8,6))
ax = fig.add_subplot(111, projection='3d')

# Axes limits
ax.set_xlim(-4,4)
ax.set_ylim(-4,4)
ax.set_zlim(-1,4)

# Plot column vectors
ax.quiver(0,0,0, *a1, color='blue', label='a1')
ax.quiver(0,0,0, *a2, color='green', label='a2')
ax.quiver(0,0,0, *a3, color='red', label='a3')
ax.quiver(0,0,0, *b,  color='purple', label='b')

# Plot scaled solution vectors
ax.quiver(0,0,0, *v1, color='blue', alpha=0.5)
ax.quiver(0,0,0, *v2, color='green', alpha=0.5)
ax.quiver(0,0,0, *v3, color='red', alpha=0.5)

# Polygonal chain to show addition step by step
ax.plot([0,p1[0]],[0,p1[1]],[0,p1[2]],'k--')
ax.plot([p1[0],p2[0]],[p1[1],p2[1]],[p1[2],p2[2]],'k--')
ax.plot([p2[0],p3[0]],[p2[1],p3[1]],[p2[2],p3[2]],'k--')

\end{lstlisting}
\end{frame}

\begin{frame}[fragile]{Python Code}
\begin{lstlisting}[language=Python]

# Labels
ax.text(*a1, "a1", color='blue')
ax.text(*a2, "a2", color='green')
ax.text(*a3, "a3", color='red')
ax.text(*b, "b", color='purple')
ax.text(*v1, "2a1", color='blue')
ax.text(*v2, "1/2 a2", color='green')
ax.text(*v3, "2/3 a3", color='red')

# Axes labels and title
ax.set_xlabel("x")
ax.set_ylabel("y")
ax.set_zlabel("z")
ax.legend()
ax.set_title("3D visualization of 2a1 + (1/2)a2 + (2/3)a3 = b")

plt.show()
\end{lstlisting}
\end{frame}

%------------------------------------

\begin{frame}[fragile]{C Code}
\begin{lstlisting}[language=C]
#include <stdio.h>

#define N 3  // 3x3 matrix

// Function to calculate determinant of a 3x3 matrix
float determinant(float A[3][3]) {
    float det;
    det = A[0][0]*(A[1][1]*A[2][2] - A[1][2]*A[2][1])
        - A[0][1]*(A[1][0]*A[2][2] - A[1][2]*A[2][0])
        + A[0][2]*(A[1][0]*A[2][1] - A[1][1]*A[2][0]);
    return det;
}

// Function to find adjoint of 3x3 matrix
void adjoint(float A[3][3], float adj[3][3]) {
    adj[0][0] =   (A[1][1]*A[2][2] - A[1][2]*A[2][1]);
    adj[0][1] = - (A[0][1]*A[2][2] - A[0][2]*A[2][1]);
    adj[0][2] =   (A[0][1]*A[1][2] - A[0][2]*A[1][1]);

    adj[1][0] = - (A[1][0]*A[2][2] - A[1][2]*A[2][0]);
    adj[1][1] =   (A[0][0]*A[2][2] - A[0][2]*A[2][0]);
    adj[1][2] = - (A[0][0]*A[1][2] - A[0][2]*A[1][0]);

    adj[2][0] =   (A[1][0]*A[2][1] - A[1][1]*A[2][0]);
    adj[2][1] = - (A[0][0]*A[2][1] - A[0][1]*A[2][0]);
    adj[2][2] =   (A[0][0]*A[1][1] - A[0][1]*A[1][0]);

    // Transpose of cofactor matrix = adj(A)
    float temp;
    temp = adj[0][1]; adj[0][1] = adj[1][0]; adj[1][0] = temp;
    temp = adj[0][2]; adj[0][2] = adj[2][0]; adj[2][0] = temp;
    temp = adj[1][2]; adj[1][2] = adj[2][1]; adj[2][1] = temp;
}

// Function to find inverse of 3x3 matrix
int inverse(float A[3][3], float inv[3][3]) {
    float det = determinant(A);
    if (det == 0) {
        printf("Matrix is singular, no inverse exists.\n");
        return 0;
    }

    float adj[3][3];
    adjoint(A, adj);

    for (int i=0; i<3; i++)
        for (int j=0; j<3; j++)
            inv[i][j] = adj[i][j] / det;

    return 1;
}

// Multiply 3x3 matrix with 3x1 vector
void multiply(float A[3][3], float b[3], float x[3]) {
    for (int i=0; i<3; i++) {
        x[i] = 0;
        for (int j=0; j<3; j++)
            x[i] += A[i][j] * b[j];
    }
}

int main() {
    // Define A and b
    float A[3][3] = {
        {1, 2, -3},
        {2, 0, -3},
        {1, 2,  0}
    };
    float b[3] = {1, 2, 3};
    float invA[3][3], x[3];

    // Find inverse
    if (inverse(A, invA)) {
        // Solve x = A^{-1} b
        multiply(invA, b, x);

        printf("Solution:\n");
        printf("x = %.3f\n", x[0]);
        printf("y = %.3f\n", x[1]);
        printf("z = %.3f\n", x[2]);
    }

    return 0;
}
\end{lstlisting}
\end{frame}

\begin{frame}[fragile]{C Code}
\begin{lstlisting}[language=C]


    adj[1][0] = - (A[1][0]*A[2][2] - A[1][2]*A[2][0]);
    adj[1][1] =   (A[0][0]*A[2][2] - A[0][2]*A[2][0]);
    adj[1][2] = - (A[0][0]*A[1][2] - A[0][2]*A[1][0]);

    adj[2][0] =   (A[1][0]*A[2][1] - A[1][1]*A[2][0]);
    adj[2][1] = - (A[0][0]*A[2][1] - A[0][1]*A[2][0]);
    adj[2][2] =   (A[0][0]*A[1][1] - A[0][1]*A[1][0]);

    // Transpose of cofactor matrix = adj(A)
    float temp;
    temp = adj[0][1]; adj[0][1] = adj[1][0]; adj[1][0] = temp;
    temp = adj[0][2]; adj[0][2] = adj[2][0]; adj[2][0] = temp;
    temp = adj[1][2]; adj[1][2] = adj[2][1]; adj[2][1] = temp;
}

// Function to find inverse of 3x3 matrix
int inverse(float A[3][3], float inv[3][3]) {
    float det = determinant(A);
    if (det == 0) {
        printf("Matrix is singular, no inverse exists.\n");
        return 0;
    }

    float adj[3][3];
    adjoint(A, adj);

    for (int i=0; i<3; i++)
        for (int j=0; j<3; j++)
            inv[i][j] = adj[i][j] / det;

    return 1;
}

// Multiply 3x3 matrix with 3x1 vector
void multiply(float A[3][3], float b[3], float x[3]) {
    for (int i=0; i<3; i++) {
        x[i] = 0;
        for (int j=0; j<3; j++)
            x[i] += A[i][j] * b[j];
    }
}

int main() {
    // Define A and b
    float A[3][3] = {
        {1, 2, -3},
        {2, 0, -3},
        {1, 2,  0}
    };
    float b[3] = {1, 2, 3};
    float invA[3][3], x[3];

    // Find inverse
    if (inverse(A, invA)) {
        // Solve x = A^{-1} b
        multiply(invA, b, x);

        printf("Solution:\n");
        printf("x = %.3f\n", x[0]);
        printf("y = %.3f\n", x[1]);
        printf("z = %.3f\n", x[2]);
    }

    return 0;
}
\end{lstlisting}
\end{frame}

\begin{frame}[fragile]{C Code}
\begin{lstlisting}[language=C]
// Function to find inverse of 3x3 matrix
int inverse(float A[3][3], float inv[3][3]) {
    float det = determinant(A);
    if (det == 0) {
        printf("Matrix is singular, no inverse exists.\n");
        return 0;
    }

    float adj[3][3];
    adjoint(A, adj);

    for (int i=0; i<3; i++)
        for (int j=0; j<3; j++)
            inv[i][j] = adj[i][j] / det;

    return 1;
}

// Multiply 3x3 matrix with 3x1 vector
void multiply(float A[3][3], float b[3], float x[3]) {
    for (int i=0; i<3; i++) {
        x[i] = 0;
        for (int j=0; j<3; j++)
            x[i] += A[i][j] * b[j];
    }
}

int main() {
    // Define A and b
    float A[3][3] = {
        {1, 2, -3},
        {2, 0, -3},
        {1, 2,  0}
    };
    float b[3] = {1, 2, 3};
    float invA[3][3], x[3];

    // Find inverse
    if (inverse(A, invA)) {
        // Solve x = A^{-1} b
        multiply(invA, b, x);

        printf("Solution:\n");
        printf("x = %.3f\n", x[0]);
        printf("y = %.3f\n", x[1]);
        printf("z = %.3f\n", x[2]);
    }

    return 0;
}
\end{lstlisting}
\end{frame}

\begin{frame}[fragile]{C Code}
\begin{lstlisting}[language=C]


// Multiply 3x3 matrix with 3x1 vector
void multiply(float A[3][3], float b[3], float x[3]) {
    for (int i=0; i<3; i++) {
        x[i] = 0;
        for (int j=0; j<3; j++)
            x[i] += A[i][j] * b[j];
    }
}

int main() {
    // Define A and b
    float A[3][3] = {
        {1, 2, -3},
        {2, 0, -3},
        {1, 2,  0}
    };
    float b[3] = {1, 2, 3};
    float invA[3][3], x[3];

    // Find inverse
    if (inverse(A, invA)) {
        // Solve x = A^{-1} b
        multiply(invA, b, x);

        printf("Solution:\n");
        printf("x = %.3f\n", x[0]);
        printf("y = %.3f\n", x[1]);
        printf("z = %.3f\n", x[2]);
    }

    return 0;
}
\end{lstlisting}
\end{frame}

\begin{frame}[fragile]{C Code}
\begin{lstlisting}[language=C]

    float invA[3][3], x[3];

    // Find inverse
    if (inverse(A, invA)) {
        // Solve x = A^{-1} b
        multiply(invA, b, x);

        printf("Solution:\n");
        printf("x = %.3f\n", x[0]);
        printf("y = %.3f\n", x[1]);
        printf("z = %.3f\n", x[2]);
    }

    return 0;
}
\end{lstlisting}
\end{frame}

\begin{frame}{Plot}
    \centering
    \includegraphics[width=\columnwidth, height=0.8\textheight, keepaspectratio]{beamer/figs/fig1.png}     
\end{frame}

\end{document}
